\documentclass{article}

% Essential math packages
\usepackage{amsmath}
\usepackage{amssymb}
\usepackage{amsthm}
\usepackage{mathtools}
\usepackage{hyperref}
\usepackage{algorithm}
\usepackage{algpseudocode}
\usepackage{tikz}
\usetikzlibrary{decorations.pathreplacing,calc,arrows.meta, decorations.text}

\def\xscale{1.45}      % horizontal stretch
\def\yscale{1.25}      % vertical stretch
\def\Ysep{-6.3}        % drop of local ξ–axis
\def\BraceRaise{8pt}   % lift of braces above x–axis

\usepackage{geometry}
\usepackage{graphicx}
\usepackage{enumitem}
\usepackage{float}      
\usepackage{setspace}
\usepackage{cite}        



\setstretch{1.2}

% Page margins
\geometry{margin=1in}

\newcommand{\dif}{\mathrm{d}}
% Hyperref setup (optional, improves PDF navigation)
\hypersetup{
    colorlinks=true,
    linkcolor=blue,
    citecolor=blue,
    filecolor=magenta,
    urlcolor=cyan,
    pdftitle={The Brachistrochrone Problem: a Finite Element Approach},
    pdfauthor={Guizhong (Harry) Luo},
}

% Title information
\title{ 
  \textbf{A Strategist's Guide to Campus Dining:} \\
  \Large A Multifactor Optimization Approach \\
  \vspace{0.5em}
  \normalsize ISyE 524 Project \textit{Proposal}
}
\begin{document}


\maketitle\begin{center}
\begin{tikzpicture}
    \def\radius{0.48cm} % Adjust radius for spacing
    
    % Upper semicircle: Harry Luo (clockwise path, no draw)
    \path[decoration={
        text along path,
        text={Harry Luo},
        text align={center},
        raise=3pt},  % Raise for outside alignment
    decorate] (90:\radius) arc (90:-90:\radius);
    
    % Lower semicircle: Yufei Shen (counterclockwise, reverse for readability, no draw)
    \path[decoration={
        text along path,
        text={Yufei Shen},
        text align={center},
        raise=3pt,  % Negative for outside alignment
        reverse path},  % Reverse to make text upright
    decorate] (180:\radius) arc (90:270:\radius);
\end{tikzpicture}
\end{center}



\section{Overview}

A Chinese wisdom has been driving generations of innovators and entrepreneurs: ``Time is money; efficiency is life.'' This principle, first proclaimed during China's economic reforms in 1981, resonates deeply with the daily reality of college students navigating campus life. Indeed, time and efficiency have become the constant pursuit of students hurrying between classes, strategically timing their meals to avoid crowded dining halls, and calculating the optimal schedules across campus to maximize precious minutes between lectures.

This project is our answer to this challenge. We optimize for weekly lunch decisions given the course schedule and daily targets chosen at random. The daily targets affect our food emphasis (weights) for the day, which is a weighted combination of time cost, money cost, enjoyment, and nutrition. Satisfying the four-factor target yields an daily score, which motivates our target function: maximizing weekly score.


Real data of the transit time between lecture halls and restaurants will be gathered as time cost. For simplicity, we assume that each restaurant only produces one food option, whose nutritional value will be generated by AI.



\section{Variables(tentative)}
\begin{itemize}[left=0pt]
  \item $D=\{1,\dots,5\}$: days of the week.
  \item $R=\{1,\dots,n\}$: restaurants.
  \item $t_{d,r}$: travel time (min) on day $d$ to/from restaurant $r$(we might need to make it more precise since the starting point and the destination is tentative) .
  \item $p_r$: price at restaurant $r$.
  \item $(c_r,\,u_r,\,f_r,\,h_r)$: calories, protein, fat, carbs at $r$.
  \item $m_r$: the discount given in restaurant $r$ (could be a single constraint or just a variable in constraint(tentative right now)).
  \item $T_{r,k},\,P_{r,k},\,(c_{r,k},u_{r,k},f_{r,k},h_{r,k})$: threshold, price, and nutrients for deal $k$ at $r$.
  \item $L$: dine-in time (can be reduced if today we are in a "hurry") \item $B$: weekly budget. 
  \item $F$: fridge cap(we are not sure whether we should consider backlogging for lunch right now, but this could be a reasonable parameter), 
  \item $\sigma$: spoilage cost.(tentative to add)
  \item $L_d$: Time restriction on day $d$.
\end{itemize}

\section{Model and Method }

We'll start with a Minimum Cost Network Flow (MCNF) model focusing on time minimization, then gradually add objective variables and constraints. The core analysis will explore different weight parameter combinations, treating this as a multi-objective tradeoff problem. We'll compare two approaches: optimizing dual objectives with others as constraints versus full multi-objective optimization. Integer constraints will handle discrete choices like locations and meal sizes.

Each variable represents real student concerns—time, cost, nutrition, satisfaction, distance. Our methodology progresses from independent objectives to interconnected relationships, evolving from simplified models to realistic decision-making scenarios.





\end{document}
